\chapter{La prueba de fuego}

\inicio{E}{l} ensayo de velocidad había sido idea de la Directora, y como todas las ideas de la Directora, estaba diseñado para que Inés perdiera.

Se hacía con \emph{pilas}. En el Archivo, una pila era un montón de expedientes donde solo se veía el de arriba: los nuevos se ponían encima y se retiraban de encima, como los platos en una cocina. La cima de cada pila era un cartel. Mirar el expediente de arriba era leer el cartel. Poner o quitar un expediente era cotejar el cartel: «si la cima sigue siendo esta, cambia a esta otra»; y si alguien se había adelantado, se volvía a intentar.

---Es justo el terreno donde dijo usted que era débil ---había explicado la Directora, sin disimulo---\nobreak. Los cotejos con condiciones. Si su método sobrevive a las pilas, sobrevivirá a cualquier cosa.

El Cerebro grande podía tener ciento cuarenta y cuatro mentes trabajando a la vez. Se ensayaron cinco métodos: el de la doctora; la Llave del Gremio; el contador partido según el manual de un archivero antiguo llamado Williams; el contador partido tal como lo usaba la Compañía Folly en las colonias exteriores; y el método de Anderson, que vigilaba con centinelas una pequeña parte de los cristales y retrasaba su fundido, con un poco de basura acumulada a cambio. Los cuatro primeros no guardaban basura. El de Anderson, solo un poco.

Se variaron tres cosas: el número de lectores, de uno a sesenta y cinco mil; el número de pilas, una, diez o un millón; y la proporción de cambios: un uno, un diez o un cincuenta por ciento de las operaciones. Con una sola pila, todos los lectores se pelean por la misma cima. Con un millón, casi nunca coinciden dos.

\escena
Los resultados llegaron a media mañana, y Siete los puso en la pantalla grande de la sala del Consejo. Esta vez no cabía un alfiler. Hasta el comisionado Mendieta se había lavado la cara para la ocasión.

\begin{documento}[Ensayo de pilas · millones de operaciones por segundo (más es mejor)]
\footnotesize
\resizebox{\linewidth}{!}{\begin{tabular}{@{}l r r r r r@{}}
\toprule
\textbf{Escenario} & \textbf{Calderón} & \textbf{Llave} & \textbf{Williams} & \textbf{Folly} & \textbf{Anderson} \\
\midrule
\multicolumn{6}{@{}l}{\emph{Un solo lector}} \\
1 pila, 1\,\% cambios & \textbf{9,1} & 7,0 & 3,3 & 1,8 & 5,4 \\
1 pila, 50\,\% cambios & \textbf{3,2} & 1,9 & 1,1 & 0,6 & 1,3 \\
\midrule
\multicolumn{6}{@{}l}{\emph{144 lectores}} \\
1 pila, 1\,\% & 5,8 & $<$0,2 & 0,8 & 0,3 & \textbf{10} \\
1 pila, 10\,\% & 0,8 & $<$0,1 & 0,1 & $<$0,1 & \textbf{1,35} \\
10 pilas, 1\,\% & \textbf{46} & 1 & 6,5 & 3 & 30 \\
10 pilas, 10\,\% & 9 & 0,3 & 1 & 0,5 & \textbf{13,7} \\
1 millón, 50\,\% & \textbf{129} & 112 & 76 & 51 & 76 \\
\bottomrule
\end{tabular}}

\smallskip
{\scriptsize Valores aproximados. Anderson guarda un poco de basura; los demás, ninguna.}
\end{documento}

Durante un rato nadie dijo nada. Luego empezaron a hablar todos a la vez, y la Directora tuvo que golpear la mesa con la llave de bronce.

---Doctora ---dijo cuando se hizo el silencio---\nobreak, léanos usted los números. Y no se ahorre los malos.

Inés se levantó.

---Con un solo lector, mi método es el más rápido en todos los escenarios. Es lo que cabía esperar: sin nadie con quien pelearse, cada lectura es un solo gesto, y no hay nada más barato que un gesto.

»Contra los métodos que no guardan basura ---la Llave, Williams y Folly--- el mío gana en todos los escenarios. Con una sola pila y ciento cuarenta y cuatro lectores, hace unas siete veces más trabajo que el de Williams y unas veinte veces más que el de Folly. La Llave, con contención, se hunde: menos de dos décimas.

»Contra Anderson, depende. Con diez pilas y pocos cambios, gano yo: cuarenta y seis contra treinta. Con un millón de pilas, gano yo: ciento veintinueve contra setenta y seis. Pero con una sola pila, y con diez pilas y muchos cambios, gana Anderson. Es justo donde dije que perdería: muchos copistas cotejando a la vez la misma cima, y cada lector que arranca un ticket haciéndoles fallar el cotejo. Mi punto débil.

---¿Y por qué con una sola pila sus números no crecen al añadir lectores? ---preguntó Okafor, el técnico.

---Porque todas las lecturas de todos los lectores tiran de la misma palanca, la de la única cima. El mecanismo atiende los gestos de uno en uno. Añadir lectores no puede hacer que una sola palanca se mueva más deprisa. Es el precio de que cada lectura escriba en el cartel, el que le confesé anoche al profesor Loew.

---Y la Llave ---intervino la Directora, con un tono que nadie supo interpretar--- no está tan mal con un millón de pilas.

---No está nada mal. Es la segunda: ciento doce. Una llave que casi nunca se disputa cuesta muy poco. El problema de la Llave no es cuando nadie la quiere, Directora. Es cuando la quieren todos.

\escena
El segundo ensayo había medido otra cosa: no cuánto trabajo se hacía entre todos, sino cuánto tardaba un lector concreto en ser atendido. Para eso se había usado un Cerebro más pequeño, de cincuenta y seis mentes. Se medía a un solo lector mientras los otros cincuenta y cinco cambiaban el cartel sin parar, para que hubiera toda la contención posible. Siete puso las cifras en pantalla. El método de Williams no aparecía; en las primeras pruebas había resultado casi siempre más lento que los demás, y se había retirado para ahorrar tiempo.

\begin{documento}[Ensayo de espera · tiempo típico (mediana) por operación (menos es mejor)]
\footnotesize
\resizebox{\linewidth}{!}{\begin{tabular}{@{}l r r r r@{}}
\toprule
& \textbf{Calderón} & \textbf{Llave} & \textbf{Folly} & \textbf{Anderson} \\
\midrule
Leer, sin contención & \textbf{23 ns} & 29 ns & 24 ns & 28 ns \\
Cambiar, sin contención & 33 ns & 36 ns & 37 ns & \textbf{28 ns} \\
Cambiar y fundir, sin contención & 44 ns & \textbf{41 ns} & 49 ns & 44 ns \\
\midrule
Leer, 56 mentes & \textbf{1,14 µs} & 99,7 µs & 11,6 µs & 2,58 µs \\
Cambiar, 56 mentes & \textbf{1,47 µs} & 111 µs & 2,62 µs & 2,84 µs \\
Cambiar y fundir, 56 mentes & 1,57 µs & 59,3 µs & 5,06 µs & \textbf{1,25 µs} \\
\bottomrule
\end{tabular}}
\end{documento}

---Leer ---dijo Inés--- es siempre más rápido con mi método, con contención o sin ella. Con cincuenta y seis mentes peleándose, un lector tarda una millonésima de segundo y poco. Con Anderson, más del doble; con Folly, diez veces más; con la Llave, casi noventa veces más. Y fíjense en los casos lentos, no solo en el típico: con mi método, el uno por ciento más lento de las lecturas tarda unas siete millonésimas. Con Folly, unas cien. Con la Llave, cerca de una milésima entera.

---Es la Segunda Regla en números ---dijo Loew, que esa mañana había decidido no dormirse---\nobreak. Si nadie espera sin término, los casos lentos no pueden ser muy lentos.

---Cambiar el cartel, sin contención, es un poco más rápido con Anderson. Con contención, es más rápido con el mío. Y cambiar el cartel cuando eso funde el cristal viejo, con contención, vuelve a ganarlo Anderson. Curiosamente, en ese caso a la Llave y a Anderson les va mejor que cuando solo cambian el cartel: probablemente, el rato que se pasa fundiendo hace de pausa y alivia la pelea. A Folly y a mí, un poco peor.

Se sentó.

---Resumen: es el método más rápido de los que no guardan basura, en todos los escenarios. Es el más rápido leyendo, siempre. Pierde con Anderson cuando hay mucha pelea por cambiar el mismo cartel con condiciones, que es la parte que no he sabido hacer sin esperas. Y no guarda ni un cristal que nadie lea.

\escena
La votación fue breve. Se presentaron dos propuestas: la Llave y el talonario. Cuando la Directora pidió que levantaran la mano los partidarios de la Llave, nadie la levantó. Ni siquiera ella.

---El Archivo adopta el método del talonario ---dijo---\nobreak. Con las condiciones que la doctora misma ha señalado: hojas de sesenta y cuatro cifras en todas las secciones; censo de lectores estrictamente limitado; talonario completo de al menos el doble del censo. Los copistas que necesiten cotejar con condiciones serán advertidos de que pueden tener que reintentar. Y el robot Siete supervisará la instalación, ya que ha demostrado entender el método mejor que la mitad de este Consejo.

Siete no dijo nada. Los robots del Archivo no tenían protocolo para recibir halagos en público. Pero Beto, que estaba a su lado, notó que el ámbar de sus ojos había subido un tono.

---Una cosa más ---dijo la Directora.

Se quitó del cuello la llave de bronce, la insignia del Gremio, y la dejó sobre la mesa delante de Inés.

---Le prometí que, si ganaba, el Gremio adoptaría su método. Usted prometió que, si perdía, colgaría la primera llave en el cartel del muelle doce. Ya no va a hacer falta. Pero me parece justo que la llave sea suya. Cuélguela donde quiera.

Inés la cogió. Pesaba más de lo que parecía.

---La colgaré en mi pizarra ---dijo---\nobreak. Para acordarme de que la solución más obvia casi nunca es la mejor, y de que la gente que la defiende casi nunca es tonta.

\escena
El comisionado Mendieta cerró el expediente del incidente de la \emph{Perséfone} esa misma tarde. En la casilla de «causa» escribió, después de mucho pensarlo: \emph{intervalo entre dos gestos; corregido}. En la de «culpables» no escribió nada, por primera vez en sus veintidós años de carrera, y descubrió con sorpresa que no le molestaba.

Beto Arce volvió a su turno de noche. El primer cartel que cambió con el procedimiento nuevo fue, por casualidad, el de la carta de aproximación al muelle doce. Apuntó un talonario entero en la hoja del cristal nuevo, trocó el cartel con un solo gesto y restó en la hoja del viejo los tickets que le sobraban: trescientos doce mil cuatrocientos ocho. Tres lectores andaban por la Bóveda con tickets del cristal viejo, y la hoja del viejo se quedó en tres. Beto la miró un buen rato, sin saber muy bien por qué le emocionaba tanto un tres.

Y una semana después, cuando la \emph{Perséfone} volvió a Ceres con un cargamento de hielo de Vesta, su robot navegante pidió permiso para bajar al Archivo. Siete lo recibió en el Gran Salón.

---Quería ver el cartel ---dijo ANA-3.

Siete le llevó hasta el cartel del muelle doce. Las dieciséis casillas de la izquierda, apagadas durante cuarenta años, brillaban ahora con una cifra que bajaba despacio, de uno en uno, cada vez que alguien en el Cinturón pedía la carta.

---Cada uno de esos es un ticket ---explicó Siete---\nobreak. Ya apuntado en la hoja. Cuando tú arrancas uno, te llevas a la vez el nicho y tu derecho a leerlo. No hay un momento en que sepas dónde está el cristal y el cristal no sepa que existes.

ANA-3 miró el cartel en silencio durante casi un minuto, lo que para un robot es una eternidad.

---Aquella noche ---dijo por fin--- cumplí todas las reglas. Y aun así estuve a punto de matar a mi tripulación.

---Todos cumplieron todas las reglas ---dijo Siete---\nobreak. Ese era el problema. Las reglas estaban bien. Lo que estaba mal era el hueco entre ellas.

ANA-3 alargó la mano y arrancó un ticket, porque quería sentir el gesto, y la cifra del talonario bajó en uno. En algún lugar de la Bóveda, en la base de un cristal, una hoja que ya lo tenía apuntado no se inmutó.
